\section*{NeurIPS Paper Checklist}

\begin{enumerate}

\item {\bf Claims}
    \item[] Question: Do the main claims made in the abstract and introduction accurately reflect the paper's contributions and scope?
    \item[] Answer: \answerYes{}
    \item[] Justification: The abstract and introduction state three contributions: (i) a representation-theoretic diagnosis that the Cl(3,0) geometric product realises only the L=0 and L=1 components of $1\otimes 1 = 0\oplus 1\oplus 2$ at the per-edge bilinear (Section 3.1); (ii) a constructive identification (Theorem 1, Section 3.2) showing that the geometric product combined with closed-form symmetric-traceless operations spans the Clebsch--Gordan decomposition for irrep tensor products through L=3, with explicit verification against Wigner 3j-symbols; and (iii) the CliffordSTF architecture and its empirical evaluation (Section 4): adding STF tracks raises rMD17 force cosine similarity from 0.055 to 0.551 at matched parameter budget while improving force and energy MAE, a twelve-variant controlled ablation including a readout-only control isolates the dominant effect to the L$\geq$2 substrate operating throughout message passing, and CliffordSTF outperforms a fully-configured EquiformerV2 at L=4 on OC22 IS2RE in our controlled-budget study. The scope is explicitly restricted to a controlled $\sim 10^6$-parameter protocol; we do not claim leaderboard SOTA and acknowledge this in Sections 4.1 and 5.

\item {\bf Limitations}
    \item[] Question: Does the paper discuss the limitations of the work performed by the authors?
    \item[] Answer: \answerYes{}
    \item[] Justification: Section 5 (Limitations) discusses: (i) wall-clock cost (3.39 s/step versus 0.78 s/step for EquiformerV2 at L=4, attributed to unfused PyTorch primitives); (ii) gradient-clipping requirement on QM9 and Molecule3D due to bilinear cross-track coupling; (iii) the controlled fixed-budget protocol that causes several baselines (notably MACE and ICTP) to fall short of their published accuracies, with explicit acknowledgment that our results characterise the sub-10M-parameter regime rather than the leaderboard frontier; and (iv) restriction of the directional-fidelity analysis to rMD17 small organic molecules. Section 6 additionally notes that O(3) (parity) equivariance would require explicit polar/axial channel tracking (Remark 1, Section 3.1).

\item {\bf Theory assumptions and proofs}
    \item[] Question: For each theoretical result, does the paper provide the full set of assumptions and a complete (and correct) proof?
    \item[] Answer: \answerYes{}
    \item[] Justification: Theorem 1 (Section 3.2) is stated with explicit assumptions ($u, v \in \mathbb{R}^3$ carrying the L=1 irrep of SO(3)). A proof sketch is given in the main text expanding $u_i v_j$ into symmetric-traceless, antisymmetric, and trace components and using the Hodge isomorphism $\Lambda^2 \mathbb{R}^3 \cong \mathbb{R}^3$. Full coordinate-level proofs, the higher-order couplings ($2\otimes 1, 2\otimes 2, 3\otimes 1$), the closed-form detracing coefficients, and machine-precision numerical verification against Wigner 3j-symbols are provided in Supplementary~\ref{supp:cg_proof} and~\ref{supp:cg_test}. Equivariance of all architectural operations is justified in Section 3.3 and verified empirically; classical references (Coope et al. 1965, Applequist 1989, Thorne 1980, Stone 2013, Shao et al. 2025) are cited for the underlying Cartesian--spherical equivalence.

\item {\bf Experimental result reproducibility}
    \item[] Question: Does the paper fully disclose all the information needed to reproduce the main experimental results of the paper to the extent that it affects the main claims and/or conclusions of the paper?
    \item[] Answer: \answerYes{}
    \item[] Justification: Section 4.1 specifies the unified training framework: $\sim 10^6$ parameters, Adam with learning rate $10^{-4}$, no weight decay, early stopping at patience 30, no EMA or warm-up. Per-benchmark settings (epochs, batch size, loss, seeds) are in Supplementary~\ref{supp:protocol}, and per-benchmark model architecture hyperparameters (channels, layers, $L_{\max}$, RBF) are tabulated in Supplementary~\ref{supp:model_settings}. Architectural details (embedding, dual-track message passing, cross-track coupling, update, readout) are specified in Section 3.3 and Section 3.4. Closed-form STF projector coefficients and the $8\times 8$ Cayley table specification are in Supplementary~\ref{supp:stf_coefficients}--\ref{supp:cg_proof}. The five ablation flags (Supplementary~\ref{supp:variants}) parameterise eleven of the twelve variants used in Table 3; the twelfth (\emph{STF only output}) is a structural ablation also documented in Supplementary~\ref{supp:variants}. Anonymised code is released at \url{https://anonymous.4open.science/r/Anon-CliffordSTF}.

\item {\bf Open access to data and code}
    \item[] Question: Does the paper provide open access to the data and code, with sufficient instructions to faithfully reproduce the main experimental results, as described in supplemental material?
    \item[] Answer: \answerYes{}
    \item[] Justification: The complete PyTorch implementation of CliffordSTF and every baseline, together with exact per-benchmark hyperparameters, is released at \url{https://anonymous.4open.science/r/Anon-CliffordSTF} (Reproducibility Statement, Section 6). All datasets used (rMD17, QM9, Molecule3D, OC20, OC22) are public and widely adopted; access instructions and data preprocessing scripts are included in the released repository.

\item {\bf Experimental setting/details}
    \item[] Question: Does the paper specify all the training and test details (e.g., data splits, hyperparameters, how they were chosen, type of optimizer) necessary to understand the results?
    \item[] Answer: \answerYes{}
    \item[] Justification: Data splits are specified in Section 4.1: rMD17 uses 900/100/1000 train/val/test per molecule (following the standard rMD17 protocol), QM9 and Molecule3D use standard splits, OC20 uses a 200k subset with IS2RE and S2EF splits, OC22 uses full IS2RE-Total and S2EF-Total. Optimiser, learning rate, weight decay, and early-stopping are specified in Section 4.1; per-benchmark epochs, batch sizes, loss formulations, and number of seeds are in Supplementary~\ref{supp:protocol}. Baseline architecture choices (e.g., $L_{\max} \in \{1,2,3\}$ for MACE/ICTP and $\{2,3,4\}$ for EquiformerV2 where applicable per benchmark) are specified in Section 4.1 and tabulated per benchmark in Supplementary~\ref{supp:model_settings}.

\item {\bf Experiment statistical significance}
    \item[] Question: Does the paper report error bars suitably and correctly defined or other appropriate information about the statistical significance of the experiments?
    \item[] Answer: \answerYes{}
    \item[] Justification: All numbers in Tables 1, 2, and 3 are reported as mean $\pm$ standard deviation across seeds (3 seeds for the main experiments, 6 seeds per (variant, molecule) for the ablation, computed as 3 random seeds $\times$ 2 independent instances). Section 4.1 states the reporting convention. The variability captured is seed-to-seed variation under the controlled training protocol. Standard deviations are reported as one-sigma values.

\item {\bf Experiments compute resources}
    \item[] Question: For each experiment, does the paper provide sufficient information on the computer resources (type of compute workers, memory, time of execution) needed to reproduce the experiments?
    \item[] Answer: \answerYes{}
    \item[] Justification: Wall-clock measurements (training and inference, per-step seconds, samples per second) for all baselines and CliffordSTF on rMD17 aspirin at batch size 32 on a single GPU are reported in Section 4.4 and Supplementary~\ref{supp:wallclock}, including parameter counts for every model. Aggregate compute for the full study (parameter budget $\sim 10^6$ per model, full ablation grid, multiple seeds across five benchmark families) is documented alongside the per-experiment training times in Supplementary~\ref{supp:protocol} and~\ref{supp:wallclock}.

\item {\bf Code of ethics}
    \item[] Question: Does the research conducted in the paper conform, in every respect, with the NeurIPS Code of Ethics?
    \item[] Answer: \answerYes{}
    \item[] Justification: The work develops an architectural component for machine-learned interatomic potentials using only public, widely adopted benchmarks (rMD17, QM9, Molecule3D, OC20, OC22) and involves no human subjects, personal data, or sensitive attributes. The Ethics Statement (Section 6) states this explicitly. Anonymity is preserved in the submission and in the released code repository.

\item {\bf Broader impacts}
    \item[] Question: Does the paper discuss both potential positive societal impacts and negative societal impacts of the work performed?
    \item[] Answer: \answerYes{}
    \item[] Justification: The Broader Impact statement (Section 6) discusses positive impact: by closing the L$\geq$2 expressivity gap of the geometric product without an e3nn-scale dependency stack, CliffordSTF lowers the barrier for practitioners who would benefit from equivariant interatomic potentials but lack specialised-library support, with the underlying ``two lineages are complementary'' principle generalising to other SO(3)-equivariant domains such as protein structure and fluid dynamics. The Ethics Statement notes that dual-use considerations are those shared with computational chemistry generally and are low relative to existing quantum-chemistry software already in widespread use.

\item {\bf Safeguards}
    \item[] Question: Does the paper describe safeguards that have been put in place for responsible release of data or models that have a high risk for misuse?
    \item[] Answer: \answerNA{}
    \item[] Justification: The released artefact is an interatomic-potential architecture and training code for public chemistry benchmarks. It is not a generative model, language model, or system with elevated dual-use risk; the released models predict energies and forces on small molecules and catalytic systems and pose no greater misuse risk than the public quantum-chemistry software they aim to accelerate.

\item {\bf Licenses for existing assets}
    \item[] Question: Are the creators or original owners of assets used in the paper properly credited and are the license and terms of use explicitly mentioned and properly respected?
    \item[] Answer: \answerYes{}
    \item[] Justification: All datasets (rMD17, QM9, Molecule3D, OC20, OC22) and baseline architectures (SchNet, DimeNet++, PaiNN, TorchMD-NET, ViSNet, NequIP, MACE, ICTP, EquiformerV2, GotenNet, FAENet) are cited at first use in Sections 2 and 4.1. Licences and dataset URLs are documented in Supplementary~\ref{supp:protocol} alongside the per-benchmark hyperparameters; standard public-research licences (MIT/Apache 2.0 for code, CC-BY for datasets where applicable) are respected.

\item {\bf New assets}
    \item[] Question: Are new assets introduced in the paper well documented and is the documentation provided alongside the assets?
    \item[] Answer: \answerYes{}
    \item[] Justification: The new asset is the CliffordSTF implementation (PyTorch source, ablation flag specifications, training scripts, hyperparameter configurations). It is released anonymously at \url{https://anonymous.4open.science/r/Anon-CliffordSTF} with documentation including a README, per-benchmark configuration files, and reproduction scripts. No new datasets are introduced.

\item {\bf Crowdsourcing and research with human subjects}
    \item[] Question: For crowdsourcing experiments and research with human subjects, does the paper include the full text of instructions given to participants and screenshots, if applicable, as well as details about compensation?
    \item[] Answer: \answerNA{}
    \item[] Justification: The work involves no crowdsourcing and no human subjects. All experiments use public computational-chemistry benchmarks.

\item {\bf Institutional review board (IRB) approvals or equivalent for research with human subjects}
    \item[] Question: Does the paper describe potential risks incurred by study participants, whether such risks were disclosed to the subjects, and whether Institutional Review Board (IRB) approvals were obtained?
    \item[] Answer: \answerNA{}
    \item[] Justification: The work involves no human subjects, so IRB approval is not applicable.

\item {\bf Declaration of LLM usage}
    \item[] Question: Does the paper describe the usage of LLMs if it is an important, original, or non-standard component of the core methods in this research?
    \item[] Answer: \answerNA{}
    \item[] Justification: LLMs are not a component of the core methodology, scientific rigor, or originality of this research. CliffordSTF is an SO(3)-equivariant interatomic-potential architecture based on Clifford-algebra geometric products and symmetric-traceless Cartesian tensors; no LLM is used at training, inference, or evaluation time. Per the NeurIPS LLM policy, declaration is not required.

\end{enumerate}